\section{Multiplicativity and a growing calibration budget}\label{sec:multiplicative}

\subsection{The class and the quantitative theorem}

Restore the actual leading phase
\(u_1=\ee^{\ii\Theta(T,a)}\). Let
\(\omega:\N\to\C\) be completely multiplicative, with
\(|\omega(n)|=1\) and \(\omega(1)=1\). Define
\begin{equation}\label{eq:qomega}
 q_\omega(w)=\frac12\sum_{n\le N}d_n
 \left[u_1\omega(n)\ee^{\ii\lambda_nw}
        +\overline{u_1\omega(n)}\ee^{-\ii\lambda_nw}\right]+K(w).
\end{equation}
The actual phase choice is \(\omega_{\can}(n)=n^{-\ii\nu}\).
For arbitrary \(\omega\), all identities
\[
 u_{mn}u_1=u_mu_n,\qquad u_n=u_1\omega(n),
\]
hold, including every retained product \(mn\le N\). The heat weights
are kept exactly; they are not multiplicative:
\begin{equation}\label{eq:weightnonmult}
 \frac{d_{mn}}{d_md_n}
       =\exp\left(\frac{c\log m\log n}{2L}\right).
\end{equation}

Let \(\cQ\) be any set of prescribed primes below \(N/2\), with arbitrary
assigned unit phases. Its cost is
\begin{equation}\label{eq:CQ}
 C_{\cQ}=\prod_{p\in\cQ}(1-p^{-1/2})^{-1}.
\end{equation}
Set
\[
 b_L=\ee^{-(1/4+c/16)L},\qquad
 M=\left\lfloor\frac{N}{40\log N}\right\rfloor,\qquad
 W=M b_L.
\]
The letter \(W\) in this section is a positive population scale; it is
distinct from the three-height quantity \(W_T(a)\).
Eventually
\begin{equation}\label{eq:Wpopulation}
 \frac{\ee^{\delta L}}{80L}\le W\le\frac{\ee^{\delta L}}{10L},
 \qquad \delta=\frac14-\frac c{16}.
\end{equation}
The useful dimensionless error budget is
\begin{equation}\label{eq:calbudget}
 \mathcal E_{\cQ}=\frac{C_{\cQ}L^4}{W}.
\end{equation}

Define fixed positive constants
\begin{align}
 a_0&=\sqrt{17/12}-1,\qquad
 \ell_1=\tfrac12\log(5/3),\qquad \ell_2=\tfrac12\log(3/2),\\
 h_2&=a_0(\ell_1^2-\ell_2^2),\qquad
 h_4=a_0(\ell_1^4-\ell_2^4),\\
 \Gamma_\beta&=\frac{\beta^6h_2h_4}{24},\qquad
 \kappa_\beta=\frac{\Gamma_\beta}{12800}.
 \label{eq:constants}
\end{align}

\begin{theorem}[Uniform multiplicative calibration obstruction]
\label{thm:mult}
Fix \(\beta>0\). Suppose the prescribed prime sets satisfy
\(\mathcal E_{\cQ}\to0\) uniformly in \eqref{eq:band}. At every
sufficiently large center there is a completely multiplicative
\(\omega\) attaining all prescribed phases such that the full
\eqref{eq:qomega}, with the actual \(u_1,d_n,K\), satisfies
\begin{equation}\label{eq:multnegative}
 q_\omega(0)=q_\omega'(0)=0,\qquad
 \cD_a(q_\omega)\le
 -\kappa_\beta\frac{\ee^{(1/2-c/8)L}}{L^8}.
\end{equation}
The conclusion holds for any prescribed phases, not only canonical
ones. Its onset is finite and uninstantiated.
\end{theorem}

The proof combines three ingredients: isolated primes that can be
adjusted independently without breaking multiplicativity; a bounded
complete complement obtained by finite character orthogonality; and
a uniform two-parameter correction of the event equations.
The general isolated-prime first-jet mechanism predates RSCT
\cite{collision}. The extension here is its negative signed margin,
complete correction, explicit calibration budget, and simultaneous
fourth-derivative compatibility in Section~\ref{sec:mirror}.

\subsection{A reserved prime population}

Use the intervals
\[
 I_1=[11/20,3/5],\quad I_2=[7/10,3/4],\quad
 I_3=[17/20,9/10].
\]
The prime number theorem \cite{dlmfpnt} gives
\[
 \#\{p:p/N\in I_j\}=(1/20+o(1))N/\log N.
\]
Select \(M\) primes from each interval and arrange them in \(M\)
triples \(p_{j,1}<p_{j,2}<p_{j,3}\). Write \(\cR\) for their union.
Every reserved prime exceeds \(N/2\); among \(n\le N\), its only
multiple is itself. Changing its phase therefore changes no
complementary term. In particular \(\cR\cap\cQ=\varnothing\).

For a reserved prime write \(v=L/2-\log p\). Direct substitution gives
\begin{equation}\label{eq:primeweights}
 d_p=b_L\ee^{v/2+cv^2/(4L)},\qquad \lambda_p=v/(2L).
\end{equation}
Thus \(d_p/b_L\) and \(L\lambda_p\) lie in fixed positive compact
intervals, separated between the three groups. The floor in \(M\)
and \(N\sim\ee^{L/2}\) give \eqref{eq:Wpopulation}. For instance,
\(M\ge N/(80\log N)\), \(N\ge\ee^{L/2}/2\), and
\(\log N\le L/2\) give its lower bound.
Only primes in fixed relative intervals have been used; no
short-interval theorem or effective numerical onset is invoked.

\subsection{The complete complement and its dependence on calibration}

Keep the prescribed phases on \(\cQ\). Assign independent uniform circle
phases to every other nonreserved prime at most \(N\). Let \(B(w)\)
be the complete sum in \eqref{eq:qomega} over \(n\notin\cR\), excluding
only the separate correction \(K\). It includes \(n=1\), every
unreserved prime, and every composite through \(N\).

\begin{lemma}[Finite orthogonality with prescribed primes]\label{lem:complement}
For any one of the observations
\[
 B^{(j)}(0)\quad(0\le j\le4),\qquad B(\ii\beta),\quad B'(\ii\beta),
\]
the mean square over the free prime phases is at most
\(9\ee^{\beta/2}C_{\cQ}^2\).
There is consequently one choice of the complementary phases for which
the sum of all seven squared observations is at most
\(63\ee^{\beta/2}C_{\cQ}^2\). If only orders zero and one at zero
and \(\ii\beta\) are requested, the corresponding constant is
\(36\ee^{\beta/2}C_{\cQ}^2\).
\end{lemma}
\begin{proof}
The elementary weight bounds are
\begin{equation}\label{eq:weightbounds}
 d_n^2\le n^{-1-c/4}\le n^{-9/8},\qquad
 \sum_{n\le N}d_n^2\le9,\qquad d_{rm}\le r^{-1/2}d_m
 \quad(rm\le N).
\end{equation}
The last follows by integrating the logarithmic slope bound
\(\dd\log d(x)/\dd\log x\le-1/2\).
Every \(n\) has a unique factorization \(n=rm\), with \(r\)
supported on \(\cQ\) and \(m\) free of primes in \(\cQ\).
Group the terms in either arm by the character \(\omega(m)\) of the
free prime variables. Distinct \(m\)'s give distinct characters.
No reserved prime enters a complementary character because its only
retained multiple was removed.

For each observation the absolute value of a character coefficient is
at most
\[
 \ee^{\beta/4}d_m
       \sum_{r:\,\cQ\text{-supported}}r^{-1/2}
 =\ee^{\beta/4}C_{\cQ}d_m;
\]
all fixed derivative factors obey \(\lambda_n^j\le1\).
Integration of \(\ee^{\ii k\theta}\) annihilates unequal characters,
giving the bound by \eqref{eq:weightbounds}. The deterministic
\(m=1\) group is included. The reflected arm has the same estimate;
\(|(A_++A_-)/2|^2\le(|A_+|^2+|A_-|^2)/2\) pays their cross term.
Sum the expectations and select a point no larger than the mean.
\end{proof}

This is finite integration on an auxiliary phase space. It is neither
a numerical experiment nor an average over actual heights. In
particular, the selected complement has controlled values at the
listed observations; we have not asserted a uniform disk bound for it.

\subsection{Negative blocks of isolated primes}

For each reserved triple apply Lemma~\ref{lem:triple} with
\[
 b_i=d_{p_{j,i}},\qquad r_i=\lambda_{p_{j,i}}.
\]
Its triangle condition holds with a uniform margin:
\[
 b_1/b_3\le(p_{j,3}/p_{j,1})^{3/4}\le(18/11)^{3/4}<2,
 \qquad b_2>b_3.
\]
After division by \(b_L\) and multiplication of frequencies by \(L\),
the block parameters form a compact nondegenerate family. Thus
\(A_2,A_3,v,\rho v\asymp b_L\).
Also
\[
 A_2\ge b_1-b_3\ge a_0b_L,\qquad
 Lr_1\ge\ell_1,\quad Lr_2\le\ell_2
\]
for large \(N\). For the last upper bound,
\(\ee^{L/2}/N\le21/20\) when \(N\ge20\).

Let \(p\) be the sum of these blocks. The positive power-series
coefficients in \eqref{eq:tripleXY} give
\begin{equation}\label{eq:primeXG}
 X\ge\frac{\beta^2h_2}{2}\frac W{L^2},\qquad
 \beta X'-2X\ge\frac{\beta^4h_4}{12}\frac W{L^4}.
\end{equation}
The first bound uses the quadratic cosh term; the second uses the
quartic term, whose factor under \(\beta\partial_\beta-2\) is two.
All blocks have aligned signs. Hence
\begin{equation}\label{eq:primeD}
 \cD_2(p)\le-\Gamma_\beta W^2/L^6,\qquad
 |p(\ii\beta)|+|p'(\ii\beta)|=O_\beta(W/L^2).
\end{equation}
At zero the same block algebra gives the derivatives needed later:
\begin{equation}\label{eq:primejets}
 p''(0)\le-h_2W/L^2,\qquad
 p^{(4)}(0)\ge h_4W/L^4.
\end{equation}
Indeed \(p''(0)=-\sum r_i^2x_i\) and
\(p^{(4)}(0)=\sum r_i^4x_i\); the \(A_2\) terms alone provide these
two bounds, and the \(A_3\) terms have the same signs.

\subsection{Uniform repair of the exact event}

In each reserved triple shift the phase of its second member by the
same \(t\), and its third by the same \(s\). The normalized event map is
\[
 F(t,s)=\left(W^{-1}\sum_{\cR}x_n(t,s),\
                   LW^{-1}\sum_{\cR}\lambda_n y_n(t,s)\right).
\]
At the block configuration \(F(0,0)=0\), and
\begin{equation}\label{eq:repairJacobian}
 DF(0)=
 \frac1W\begin{pmatrix}
 -\sum y_2&-\sum y_3\\
 L\sum r_2x_2&L\sum r_3x_3
 \end{pmatrix}
 =\begin{pmatrix}-A_J&B_J\\-C_J&-D_J\end{pmatrix},
 \qquad A_J,B_J,C_J,D_J\asymp1.
\end{equation}
The signs come from \(x_2,x_3<0\) and \(y_2>0>y_3\).
Thus \(\det DF=A_JD_J+B_JC_J\) is uniformly positive.
The inverse and the second derivatives on a fixed small parameter disk
are uniformly bounded.

Use Lemma~\ref{lem:complement} to fix a complementary phase choice.
Solve
\begin{equation}\label{eq:repairtarget}
 F(t,s)=\left(-\frac{B(0)+K(0)}W,\
                    \frac{L(B'(0)+K'(0))}W\right).
\end{equation}
The target has norm \(O_\beta(C_{\cQ}L/W)\), which tends to zero under
\eqref{eq:calbudget}. A quantitative inverse-function argument gives
\begin{equation}\label{eq:repairsize}
 |t|+|s|=O_\beta(C_{\cQ}L/W).
\end{equation}
Explicitly, apply the inverse of \eqref{eq:repairJacobian} to the
target and to the quadratic Taylor remainder. Uniformly bounded
second derivatives make this a contraction on a ball of radius a
fixed multiple of the target norm, once that norm is small.

Assign the final prime value \(\omega(p)=\ee^{\ii\theta_p}/u_1\)
at each reserved prime, keep every complementary prime value, and
extend completely multiplicatively. Isolation ensures that no
complementary coefficient changes. If \(\Delta p\) is the reserved
repair, the complete field
\[
 q_*=B+p+\Delta p+K
\]
satisfies \(q_*(0)=q_*'(0)=0\) exactly.

The repair must be estimated after the event equations have been used.
Since \(p=p^\circ\) and \(q_*=q_*^\circ\),
\[
 q_*=p+E,\qquad E=B^\circ+K^\circ+(\Delta p)^\circ.
\]
For a changed mode the subtracted value and derivative on a fixed disk
are \(O_\beta(d_n(|t|+|s|)\lambda_n^2)\).
Equations~\eqref{eq:primeweights} and~\eqref{eq:repairsize} therefore give
\[
 |(\Delta p)^\circ(\ii\beta)|
       +|((\Delta p)^\circ)'(\ii\beta)|
       =O_\beta(C_{\cQ}/L).
\]
The complete background observations from Lemma~\ref{lem:complement}
and the correction bound give
\begin{equation}\label{eq:Ebound}
 |E(\ii\beta)|+|E'(\ii\beta)|=O_\beta(C_{\cQ}).
\end{equation}
Raw repair values could be \(O(C_{\cQ}L)\); event subtraction is what
avoids that unnecessary loss.

Expanding the full quadratic expression, including background and
correction cross terms and \(\mathfrak d_a-2\), now yields
\begin{equation}\label{eq:fullprimeerror}
 \cD_a(q_*)\le-\Gamma_\beta\frac{W^2}{L^6}
 +O_\beta\left(\frac{C_{\cQ}W}{L^2}+C_{\cQ}^2
                         +a^{-2}\frac{W^2}{L^6}\right).
\end{equation}
The relative errors are respectively
\(O_\beta(\mathcal E_{\cQ})\),
\(O_\beta(\mathcal E_{\cQ}^2/L^2)\), and \(O_\beta(a^{-2})\).
After they total at most \(\Gamma_\beta/2\), the lower bound in
\eqref{eq:Wpopulation} proves
\[
 \cD_a(q_*)\le-\tfrac12\Gamma_\beta W^2/L^6
 \le-\frac{\Gamma_\beta}{12800}
                         \frac{\ee^{(1/2-c/8)L}}{L^8}.
\]
This completes the proof of Theorem~\ref{thm:mult}.
The small physical approximation error from
Proposition~\ref{prop:Derror} is not assigned to this artificial field;
no twisted physical approximation theorem has been asserted.

\subsection{Fixed and growing actual-prime calibration}

For a fixed finite set \(\mathcal P_0\) and a fixed prime
\(q_0\notin\mathcal P_0\), prescribe
\[
 \omega(p)=p^{-\ii\nu}\ (p\in\mathcal P_0),\qquad
 \omega(q_0)=-q_0^{-\ii\nu}.
\]
Here \(C_{\cQ}\) is fixed and \(\mathcal E_{\cQ}\to0\).
Theorem~\ref{thm:mult} recovers the fixed-calibration result, with the
explicit same-height defect
\begin{equation}\label{eq:phasedefect}
 |u_{q_0}-u_{q_0}^{\can}|=2.
\end{equation}

The displayed budget also permits a growing set.
\begin{corollary}[A sufficient growing calibration range]\label{cor:growing}
Fix \(0<\gamma<9/256\), and put
\begin{equation}\label{eq:Y}
 Y=\gamma L^2(\log L)^2.
\end{equation}
For every sufficiently large center, all actual prime phases
\(\omega(p)=p^{-\ii\nu}\) for \(p\le Y\) can be retained in
Theorem~\ref{thm:mult}. In addition one can choose a prime
\(q_0\in(Y,2Y]\) with the opposite actual phase, so
\eqref{eq:phasedefect} holds.
\end{corollary}
\begin{proof}
Partial summation and the prime number theorem give
\[
 \sum_{p\le Y}p^{-1/2}
 =\frac{\pi(Y)}{\sqrt Y}
       +\frac12\int_2^Y\frac{\pi(t)}{t^{3/2}}\dd t
 \sim\frac{2\sqrt Y}{\log Y}.
\]
Expanding \(-\log(1-p^{-1/2})\), the terms of order at least two
total \(O(\sum_{p\le Y}p^{-1})=O(\log Y)\), already small enough here.
Thus
\[
 \log\prod_{p\le Y}(1-p^{-1/2})^{-1}
 =\frac{2\sqrt Y}{\log Y}(1+o(1))
 =(\sqrt\gamma+o(1))L.
\]
Since \(\delta\ge3/16\), \eqref{eq:Wpopulation} implies
\[
 \mathcal E_{\cQ}\le80L^5
       \exp[-(3/16-\sqrt\gamma+o(1))L]\longrightarrow0.
\]
The prime number theorem supplies a \(q_0\in(Y,2Y]\); it is below
\(N/2\) for large \(L\). Its additional cost
\((1-q_0^{-1/2})^{-1}\) does not change the budget.
\end{proof}

For example \(\gamma=1/64\) calibrates
\((1+o(1))L^2\log L/128\) primes. Complete multiplicativity then fixes
every retained \(Y\)-smooth integer at its actual phase.
The coefficient \(9/256\) is sufficient for this proof, not a claimed
sharp boundary. The explicit defect above proves failure of the actual
same-height phase law. Realizing a similar phase pattern at another
height would change \(d_n,N,u_1,K\) and the time-height relation; no
such realization has been proved.
